%% ============================================================
%%  Chapter 7 — Short-Term Memory & Working Context
%%  Production-Grade Agentic AI: LangGraph + Python Frameworks
%% ============================================================

\chapterquote{The model does not remember. You decide what it sees.}{anon}

The production failure this chapter prevents is context amnesia under
load: the condition where a long-horizon agent loop silently forgets
that a user said ``never use code blocks'' three turns ago, or truncates
the only sentence that contained a critical user preference, because
nobody managed what was placed in the context window between turns.
This is not a token-overflow error --- no exception is raised. The
model's input is technically valid. The agent's behaviour simply
degrades, and the degradation is invisible until a user complains or an
evaluation catches it. In a customer-facing agent running hundreds of
sessions per hour, that silent degradation is a support cost, a churn
driver, and an on-call incident waiting to happen.

The solution operates in three layers. The first layer is an explicit
\emph{working buffer}: a typed message list backed by LangGraph's
\code{MessagesState} and the \code{add\_messages} reducer, which gives
the graph a single source of truth for what the model has said and seen
in the current session. The second layer is a \emph{budget enforcer}:
a \code{trim\_messages} function that evicts the oldest content before
every model call, with invariants that protect the system message and
the current user turn from eviction. The third layer is a
\emph{compression pipeline}: a summarisation node that converts
accumulated history into a compact prose summary stored in the
\code{HISTORY} slot, and a \code{MemoryInjectionNode} that formats and
injects both that summary and any retrieved memories through a six-step
pipeline governed by a four-question policy function.

The sections are ordered by dependency. Section~\ref{sec:messages-state}
builds the working buffer first, because trimming and summarisation both
operate on that buffer and cannot be explained without it.
Section~\ref{sec:trim-messages} adds the budget enforcer, because
trimming is the first-line defence before summarisation is needed.
Section~\ref{sec:summarization-node} adds the summarisation node, which
produces the compressed \code{history} string consumed by the injection
pipeline. Section~\ref{sec:injection-pipeline} assembles the six-step
injection pipeline into a \code{MemoryInjectionNode}.
Section~\ref{sec:compact-format} formalises the compact memory text
format that the pipeline emits, and explains why raw record injection is
an anti-pattern. Section~\ref{sec:injection-policy} adds the
four-question policy guard that decides, before any injection, whether
the pipeline should run at all.

This chapter introduces seven domain terms that appear throughout the
rest of the book. Table~\ref{tab:memory-taxonomy} maps each term to its
scope before any implementation detail.

\begin{table}[ht]
\centering
\caption{Short-term memory terms, scope, and lifetime introduced in Chapter~7}
\label{tab:memory-taxonomy}
\begin{tabular}{llll}
\toprule
\textbf{Term} & \textbf{Scope} & \textbf{Lifetime} & \textbf{Owner} \\
\midrule
\code{MessagesState}      & In-context turn list  & Current session     & LangGraph \\
\code{add\_messages}      & Append reducer        & Per graph write     & LangGraph \\
\code{trim\_messages()}   & Budget enforcer       & Per model call      & \S\ref{sec:trim-messages} \\
\code{summarization\_node} & History compressor   & Triggered on N turns & \S\ref{sec:summarization-node} \\
\code{MemoryInjectionNode} & Six-step pipeline    & Pre-model call      & \S\ref{sec:injection-pipeline} \\
\code{MemoryType}         & Label for memory kind & Permanent in record & \S\ref{sec:compact-format} \\
\code{should\_inject\_memory} & Policy gate       & Pre-inject          & \S\ref{sec:injection-policy} \\
\midrule
Short-term memory  & \multicolumn{3}{l}{\code{MessagesState} list; lives in-context; lost on session end} \\
Long-term memory   & \multicolumn{3}{l}{Persisted by mem0; retrieved per query; introduced in Chapter~8} \\
\bottomrule
\end{tabular}
\end{table}

The source layout for this chapter is:

\begin{verbatim}
src/
  agent/
    state/
      messages_state_short_term.py   # §7.1
      trim_messages.py               # §7.2
    memory/
      summarization_node.py          # §7.3
      memory_injection_node.py       # §7.4
      format_memory_block.py         # §7.5
      injection_policy.py            # §7.6
\end{verbatim}

%% ------------------------------------------------------------
\section{The Working Buffer Has One Owner}
\label{sec:messages-state}
%% ------------------------------------------------------------

The fundamental problem with ad-hoc message management is that no
single object in the graph owns the list of what has been said. When
ownership is diffuse --- the model node appends messages, the tool node
appends messages, the retry handler occasionally removes the last
message --- the list becomes inconsistent. Two nodes may write
different views of the same turn, a retry may silently duplicate a
message, or a tool response may appear before the user turn that
triggered it. The model sees this inconsistency as contradictory
context; the output degrades in ways that are almost impossible to
reproduce because the inconsistency depends on concurrency timing and
retry history.\looseness=-1

LangGraph's \code{MessagesState} and \code{add\_messages} reducer solve
this by making the message list a reducer-owned, append-only structure.
Every node that writes to \code{state[``messages'']} does so through
\code{add\_messages}, which handles deduplication by message ID and
ensures that updates from parallel branches are merged without
conflict. No node ever replaces the entire list; every node declares
only the messages it wants to add or update. The graph runtime handles
the merge.

The boundary between short-term and long-term memory is structural, not
conceptual. \code{MessagesState} holds only what the model needs to
reason about the \emph{current session's} turns. It is in-context: every
message in the list occupies tokens on every model call. It is
session-scoped: when the session ends, the list is discarded unless
explicitly checkpointed. Long-term memory --- persistent, cross-session
facts about the user or the world --- is outside this list. Chapter~8
introduces mem0 as the store for that category of memory and shows how
its retrieval results flow into the injection pipeline built in
\S\ref{sec:injection-pipeline}. Nothing from mem0 ever lives directly
in \code{MessagesState}; it is formatted and injected as a
\code{SystemMessage} block at the injection step, keeping the turn list
semantically clean.

The \code{AgentState} TypedDict wraps \code{MessagesState} together with
the supplementary fields the rest of this chapter writes into. The
\code{history} field is the compressed prose summary produced by the
summarisation node (\S\ref{sec:summarization-node}). The \code{meta}
field is a free-form dict for per-invocation metadata such as
\code{user\_id} and \code{session\_start}, which the injection policy
(\S\ref{sec:injection-policy}) needs but which do not belong in the
message list.

\begin{lstlisting}[language=Python,
  caption={AgentState TypedDict, SYSTEM\_PROMPT bootstrap, and model\_node using add\_messages},
  label={lst:messages-state}]
# src/agent/state/messages_state_short_term.py
from __future__ import annotations
from typing import Any
from typing_extensions import TypedDict, Annotated
from langchain_core.messages import SystemMessage, BaseMessage
from langgraph.graph.message import add_messages
from langchain_core.runnables import Runnable

SYSTEM_PROMPT = (
    "You are a helpful production assistant. "
    "Follow user instructions precisely."
)

class AgentState(TypedDict):
    messages: Annotated[list[BaseMessage], add_messages]
    history:  str | None
    meta:     dict[str, Any]

def bootstrap_state() -> AgentState:
    return AgentState(
        messages=[SystemMessage(content=SYSTEM_PROMPT)],
        history=None,
        meta={},
    )

def model_node(
    state: AgentState,
    llm:   Runnable,
) -> dict:
    response = llm.invoke(state["messages"])
    return {"messages": [response]}
\end{lstlisting}

Two decisions in Listing~\ref{lst:messages-state} deserve explanation.

First, \code{messages} is declared as \code{Annotated[list[BaseMessage],
add\_messages]} rather than a plain \code{list}. The \code{Annotated}
wrapper tells LangGraph's graph runtime which reducer to use when two
branches both write to \code{messages} in the same super-step.
Without this annotation the runtime uses last-write-wins, which silently
drops one branch's output whenever the graph forks. The annotation is
the contract between the node author and the graph runtime; omitting it
is a correctness bug, not a style issue.

Second, the system message is injected once in \code{bootstrap\_state}
and never injected again by any other node. Because \code{add\_messages}
deduplicates by message ID, a second \code{SystemMessage} with the same
content would create a duplicate entry visible to the model. Any node
that needs to extend the system context --- such as the injection node
in \S\ref{sec:injection-pipeline} --- appends a \emph{separate}
\code{SystemMessage} with a distinct ID; it never modifies or replaces
the bootstrap message.

The working buffer is now well-defined. The next problem is that it
grows without bound, and every additional message costs tokens on every
model call.

%% ------------------------------------------------------------
\section{Eviction Must Follow an Invariant, Not a Heuristic}
\label{sec:trim-messages}
%% ------------------------------------------------------------

The naive token-budget solution is to truncate the message list to the
most recent \emph{K} messages before each model call. This is simple,
predictable, and wrong in exactly the cases that matter most. It evicts
the system message when \emph{K} is small relative to the session
length. It evicts the current user turn when the turn is not literally
the last message (which happens whenever tool calls, tool responses, and
assistant intermediate messages accumulate between turns). It severs
assistant messages from the user turns that prompted them, creating a
message list where the model is answering questions it cannot see. Each
of these errors degrades model output without raising any exception.

The correct approach is an invariant-preserving eviction function that
treats the message list as a structured object rather than a sequence of
bytes. Three invariants must hold after every trim: the first
\code{SystemMessage} in the list is never evicted; the most recent
\code{HumanMessage} is never evicted; and eviction always removes
oldest-first, taking messages in user+assistant pairs wherever possible
to preserve conversational coherence. These are not preferences --- they
are correctness conditions, and the trim function asserts them on exit.

The token budget comes from outside the function: it is the
\code{history} allocation in \code{SlotBudget} (introduced in
Chapter~6), read from \code{RunnableConfig}. The trim function does not
know the model's context window; it knows only how many tokens the
\code{HISTORY} slot has been allocated. This separation means a test
can pass a 200-token budget and verify deterministic eviction on a
three-turn fixture without instantiating an LLM.

\begin{warningbox}
Never evict messages individually from the middle of the list. Removing
an assistant message while keeping its preceding user message leaves the
model with a question it cannot answer and an answer it cannot attribute.
Always drop user+assistant pairs as an atomic unit from the oldest end.
\end{warningbox}

\begin{lstlisting}[language=Python,
  caption={trim\_messages: oldest-first eviction with system and user-turn invariants},
  label={lst:trim-messages}]
# src/agent/state/trim_messages.py
from __future__ import annotations
from typing import Callable, Sequence
from langchain_core.messages import BaseMessage, SystemMessage, HumanMessage

TokenCounter = Callable[[Sequence[BaseMessage]], int]

def trim_messages(
    messages:     list[BaseMessage],
    max_tokens:   int,
    count_tokens: TokenCounter,
) -> list[BaseMessage]:
    if not messages:
        return messages

    system_msgs = [m for m in messages if isinstance(m, SystemMessage)]
    last_human  = next(
        (m for m in reversed(messages) if isinstance(m, HumanMessage)), None
    )
    protected   = set(id(m) for m in system_msgs)
    if last_human is not None:
        protected.add(id(last_human))

    evictable = [m for m in messages if id(m) not in protected]
    result    = list(messages)

    while count_tokens(result) > max_tokens and evictable:
        # drop oldest pair
        drop_n = min(2, len(evictable))
        to_drop = set(id(m) for m in evictable[:drop_n])
        result   = [m for m in result   if id(m) not in to_drop]
        evictable = [m for m in evictable if id(m) not in to_drop]

    assert any(isinstance(m, SystemMessage) for m in result), \
        "trim_messages invariant violated: system message evicted"
    if last_human is not None:
        assert any(id(m) == id(last_human) for m in result), \
            "trim_messages invariant violated: current user turn evicted"
    return result
\end{lstlisting}

Three decisions in Listing~\ref{lst:trim-messages} deserve explanation.

First, protection is tracked by object identity (\code{id(m)}) rather
than by message type or content. A session can legitimately contain
multiple \code{HumanMessage} objects; only the most recent one is
protected. Comparing by \code{id} guarantees that exactly the object
at the tail of the list is the protected one, even if two human messages
have identical content.

Second, the function accepts a \code{TokenCounter} callable rather than
importing a specific tokeniser. This inversion is deliberate: in
production the counter wraps \code{tiktoken}; in tests it can be
\code{lambda msgs: len(msgs)} for deterministic fixture-based testing
without requiring the tiktoken binary. The same \code{TokenCounter}
type alias is reused by the injection policy in
\S\ref{sec:injection-policy} for cost-benefit estimation.

Third, the two \code{assert} statements at the end are not defensive
programming --- they are executable specifications. If a future
refactor breaks either invariant (for example, by reordering messages
before trimming), the assertion fires at the trim call site with a
precise message rather than silently corrupting the model's input. The
Chapter~9 unit-test suite targets these assertions directly.

With a budget enforcer in place, the working buffer can survive
indefinitely without overflowing. But trimming alone discards information
rather than compressing it. The next problem is how to preserve the
\emph{substance} of evicted history without keeping the verbatim text.

%% ------------------------------------------------------------
\section{Evicted History Must Be Compressed, Not Discarded}
\label{sec:summarization-node}
%% ------------------------------------------------------------

The problem with pure trimming is information loss. Oldest-first
eviction removes the turns most likely to contain the user's original
framing, their constraints, and the early decisions that all subsequent
turns assumed. A model that has lost those turns will re-ask questions
the user already answered, re-explain context the user already provided,
and occasionally contradict earlier commitments because it cannot see
them. In a support agent, that pattern triggers escalations. In a
coding agent, it triggers rework.

The solution is a summarisation node that intercepts messages before
they would be evicted and converts them into a compact prose summary
stored in \code{state[``history'']}. The summary is injected into the
\code{HISTORY} slot by the assembly node (Chapter~6,
\S\ref{sec:context-assembly-node}), giving the model a compressed
representation of the evicted turns without the token cost of the
originals. The trade-off is fidelity: a summary loses exact wording.
That trade-off is correct for most sessions, where the model needs the
\emph{substance} of earlier turns, not their verbatim text.

The node uses LangChain's \code{load\_summarise\_chain} with
\code{chain\_type=``map\_reduce''}. The map step runs a summarisation
prompt over each message individually, producing a short summary per
message; the reduce step combines those summaries into a single
paragraph. Map-reduce is chosen over \code{stuff} (which concatenates
all messages into one prompt) because the set of messages being
summarised can itself exceed the model's context window when history is
deep --- exactly the scenario where summarisation is triggered.

\begin{notebox}
Summarisation is triggered by a threshold, not by the trim function
itself. The recommended trigger is: when \code{len(messages)} exceeds
\code{2~*~keep\_last\_n} turns, move the oldest \code{len(messages) - keep\_last\_n}
turns through the summarisation chain and remove them from
\code{messages}. This keeps the live message list at a stable size and
prevents the summarisation chain itself from receiving an arbitrarily
long input.
\end{notebox}

\begin{lstlisting}[language=Python,
  caption={build\_history\_summarizer with map\_reduce config and summarization\_node},
  label={lst:summarization-node}]
# src/agent/memory/summarization_node.py
from __future__ import annotations
from langchain.chains.summarize import load_summarize_chain
from langchain_core.language_models import BaseLanguageModel
from langchain_core.documents import Document
from langchain_core.runnables import Runnable
from langchain_core.messages import BaseMessage
from src.agent.state.messages_state_short_term import AgentState

def build_history_summarizer(llm: BaseLanguageModel) -> Runnable:
    return load_summarize_chain(llm, chain_type="map_reduce")

def select_messages_for_summary(
    messages:     list[BaseMessage],
    keep_last_n:  int,
) -> tuple[list[BaseMessage], list[BaseMessage]]:
    if len(messages) <= keep_last_n:
        return [], messages
    return messages[:-keep_last_n], messages[-keep_last_n:]

def summarization_node(
    state:       AgentState,
    summarizer:  Runnable,
    keep_last_n: int = 10,
) -> dict:
    to_summarize, to_keep = select_messages_for_summary(
        state["messages"], keep_last_n
    )
    if not to_summarize:
        return {}
    docs = [Document(page_content=m.content) for m in to_summarize
            if m.content]
    result      = summarizer.invoke({"input_documents": docs})
    new_summary = result["output_text"].strip()
    existing    = state.get("history") or ""
    merged = (existing + "\n\n" + new_summary).strip() if existing else new_summary
    return {"messages": to_keep, "history": merged}
\end{lstlisting}

Three decisions in Listing~\ref{lst:summarization-node} deserve
explanation.

First, the node replaces \code{state[``messages'']} with \code{to\_keep}
--- the recent slice --- rather than computing a delta. LangGraph's
\code{add\_messages} reducer handles this correctly when the node
returns the trimmed list as the new \code{messages} value, because
\code{add\_messages} merges by message ID: messages present in
\code{to\_keep} are retained, messages absent are considered superseded.
Returning the full replacement list is safer than returning a deletion
diff, which would require the node to track per-message IDs explicitly.

Second, the node merges the new summary with \code{state[``history'']}
rather than replacing it. Over a very long session, summarisation fires
multiple times. Each invocation compresses a slice of older messages;
the accumulated \code{history} string grows incrementally. That growth
is bounded because each compression step reduces a sequence of verbose
messages to a paragraph; the \code{HISTORY} slot's token budget (from
Chapter~6's \code{SlotBudget}) constrains the total when
\code{context\_assembly\_node} renders the slot.

Third, \code{Document} wrapping is required because
\code{load\_summarise\_chain} expects a list of
\code{langchain\_core.documents.Document} objects, not raw strings or
\code{BaseMessage} objects. Filtering on \code{m.content} prevents
sending empty tool-call scaffolding messages --- which contain no
summarisable prose --- through the chain, keeping the map step's
token cost predictable.

With compressed history available in \code{state[``history'']}, the
graph now needs a principled way to decide which memories --- including
both the compressed history and any retrieved long-term facts --- should
be placed in front of the model before each call.

%% ------------------------------------------------------------
\section{Memory Injection Is a Six-Step Pipeline}
\label{sec:injection-pipeline}
%% ------------------------------------------------------------

The temptation is to write a helper function that concatenates all
available memory into a single string and prepends it to the system
prompt. That approach fails in four ways. It injects everything regardless
of relevance, burning tokens on facts the model does not need for this
turn. It mixes memories from different sources (session history, user
profile, past tasks) with no indication of their provenance, type, or
age. It produces a single monolithic blob that the model cannot
selectively attend to. And it provides no observable intermediate
state, so when the model produces a wrong output, there is no way to
determine whether the failure was in retrieval, formatting, or selection.

The correct design is a pipeline of six discrete, named steps, each
with a clear input contract and output type. Each step is independently
testable and independently loggable. The six steps are:
\code{retrieve\_raw\_memories}, which gathers candidates from
\code{state[``history'']}, mem0 (Chapter~8), and recent turns;
\code{filter\_by\_scope}, which constrains to the current user and
time window; \code{score\_and\_rank}, which orders by relevance to the
current query; \code{select\_top\_k}, which applies a hard cap;
\code{format\_memory\_block}, which converts selected memories to the
compact labelled text format defined in \S\ref{sec:compact-format}; and
\code{inject\_memory\_block}, which appends the formatted block to
\code{state[``messages'']} as a \code{SystemMessage}.

\begin{notebox}
These six step names are intentionally consistent with Chapter~8's
\code{MemoryReadNode} steps. Chapter~8 plugs mem0 retrieval into
\code{retrieve\_raw\_memories} without modifying any downstream step.
The pipeline contract --- a \code{list[RetrievedMemory]} flows from step
to step --- is the integration surface between this chapter and
Chapter~8.
\end{notebox}

The \code{RetrievedMemory} TypedDict is the data contract between steps.
It carries the memory's unique ID, its \code{MemoryType} label, the
datetime it was created, its prose summary, and its relevance score.
Every step receives and returns this typed structure, making the
pipeline's data flow explicit and type-checkable without runtime
reflection.

\begin{lstlisting}[language=Python,
  caption={RetrievedMemory TypedDict, six-step helpers, and memory\_injection\_node},
  label={lst:injection-node}]
# src/agent/memory/memory_injection_node.py
from __future__ import annotations
from datetime import datetime
from typing_extensions import TypedDict
from langchain_core.messages import SystemMessage
from src.agent.state.messages_state_short_term import AgentState
from src.agent.memory.format_memory_block import format_memory_block
from src.agent.memory.injection_policy import should_inject_memory

class RetrievedMemory(TypedDict):
    id:         str
    type:       str         # MemoryType literal
    created_at: datetime
    content:    str
    score:      float

def retrieve_raw_memories(state: AgentState) -> list[RetrievedMemory]:
    # Returns history summary + placeholder for mem0 (Chapter 8)
    memories: list[RetrievedMemory] = []
    if state.get("history"):
        memories.append(RetrievedMemory(
            id="history-summary", type="episodic",
            created_at=datetime.utcnow(), content=state["history"],
            score=1.0,
        ))
    return memories

def filter_by_scope(
    memories: list[RetrievedMemory],
    user_id:  str,
    now:      datetime,
) -> list[RetrievedMemory]:
    return [m for m in memories if m["score"] > 0.0]

def score_and_rank(
    memories: list[RetrievedMemory], query: str
) -> list[RetrievedMemory]:
    return sorted(memories, key=lambda m: m["score"], reverse=True)

def select_top_k(
    memories: list[RetrievedMemory], k: int
) -> list[RetrievedMemory]:
    return memories[:k]

def inject_memory_block(
    state: AgentState, memory_block: str
) -> dict:
    msg = SystemMessage(content=memory_block)
    return {"messages": [msg]}

def memory_injection_node(
    state:   AgentState,
    *,
    user_id: str,
    now:     datetime,
    top_k:   int = 5,
) -> dict:
    raw      = retrieve_raw_memories(state)
    scoped   = filter_by_scope(raw, user_id, now)
    ranked   = score_and_rank(scoped, query=_last_query(state))
    selected = select_top_k(ranked, top_k)
    if not should_inject_memory(state, selected, now, max_extra_tokens=512):
        return {}
    block = format_memory_block(selected)
    return inject_memory_block(state, block)

def _last_query(state: AgentState) -> str:
    from langchain_core.messages import HumanMessage
    msgs = state.get("messages", [])
    for m in reversed(msgs):
        if isinstance(m, HumanMessage):
            return m.content
    return ""
\end{lstlisting}

Three decisions in Listing~\ref{lst:injection-node} deserve explanation.

First, \code{retrieve\_raw\_memories} returns a single
\code{RetrievedMemory} for the compressed history summary rather than
parsing the summary into individual records. The history string is
already compressed prose; splitting it would require a second LLM call.
Chapter~8 extends this function to also call mem0's similarity search
and append its results to the same list, making the history summary and
retrieved long-term facts uniformly typed for all downstream steps.

Second, \code{memory\_injection\_node} calls
\code{should\_inject\_memory} \emph{after} selection but \emph{before}
formatting and injection. The policy check (\S\ref{sec:injection-policy})
needs the ranked, selected list --- not the raw list --- because its
cost-benefit question requires knowing the actual token cost of the
block that would be injected. Calling the policy on raw, unranked
memories would overestimate cost.

Third, the node returns an empty dict \code{\{\}} when the policy
decides not to inject. LangGraph treats an empty dict return as a no-op
state update, so the node remains in the graph unconditionally --- it
is never conditionally skipped at the edge level. This is preferable to
a conditional edge that bypasses the node, because it keeps the graph
topology stable and makes the injection decision observable via the
node's Langfuse span.

The pipeline produces a memory block, but the format of that block is
as important as its content. A poorly formatted injection degrades model
steering even when the content is correct.

%% ------------------------------------------------------------
\section{Raw Records Are an Anti-Pattern; Labels Are Mandatory}
\label{sec:compact-format}
%% ------------------------------------------------------------

The most common injection mistake is appending a Python \code{dict} or
JSON object directly into the prompt: \verb|{"type": "preference",|
\verb|"content": "user prefers dark mode", "created_at": "2026-01-15"}|.
That format fails for three reasons. It is verbose: a 12-word user
preference is wrapped in 60 characters of punctuation and field names.
It leaks implementation: the model sees internal field names, storage
timestamps in ISO format, and UUIDs that mean nothing in the context of
a reasoning task. And it degrades model steering: the model's
instruction-following is calibrated to prose, not to JSON-in-the-middle-of-prose.
A model that receives JSON fragments in its system message attends to them
less reliably than a model that receives a labelled, natural-language
sentence.

The mandatory format is a compact labelled line: \verb|Memory [type date]: summary|.
The label provides the model with two signals at low token cost: the
\emph{type} tells it what kind of fact this is (a past event, a standing
preference, or an active task), and the \emph{date} tells it how recent
the fact is, allowing the model to weight recency in its reasoning
without being given a raw timestamp to parse. The summary is always
a complete sentence or short paragraph --- never a field value
extracted from a dict. The date format is \code{YYYY-MM-DD}: unambiguous,
compact, and interpretable by the model without additional instructions.

\begin{lstlisting}[language=Python,
  caption={MemoryType, format\_single\_memory, and format\_memory\_block},
  label={lst:format-memory}]
# src/agent/memory/format_memory_block.py
from __future__ import annotations
from datetime import datetime
from typing import Literal

MemoryType = Literal["episodic", "profile", "preference", "task"]

def format_single_memory(
    mt:         MemoryType,
    created_at: datetime,
    summary:    str,
) -> str:
    date_str = created_at.strftime("%Y-%m-%d")
    return f"Memory [{mt} {date_str}]: {summary}"

def format_memory_block(memories: list) -> str:
    """Accept list[RetrievedMemory]; return multi-line memory block."""
    lines = []
    for m in sorted(memories, key=lambda x: x["created_at"]):
        lines.append(
            format_single_memory(
                mt=m["type"],
                created_at=m["created_at"],
                summary=m["content"],
            )
        )
    return "\n".join(lines)

def make_memory_system_message(memory_block: str):
    from langchain_core.messages import SystemMessage
    return SystemMessage(content=memory_block)
\end{lstlisting}

Two decisions in Listing~\ref{lst:format-memory} deserve explanation.

First, \code{format\_memory\_block} sorts memories by
\code{created\_at} ascending --- oldest first --- rather than by score
descending. When the model reads a block of memory lines, temporal
order provides a narrative coherence that relevance order does not: the
user established a preference, then modified it, then acted on it; that
sequence is meaningful. Within the block, all selected memories have
already passed the \code{select\_top\_k} filter, so they are all
considered relevant; ordering them by time is more useful than ordering
them by a score the model never sees.

Second, the four \code{MemoryType} literals are not open-ended strings.
Constraining the type vocabulary prevents two common errors: a downstream
step that branches on memory type (as Chapter~8's write node does) will
fail with a clear \code{mypy} error if an unrecognised type is passed,
rather than silently reaching a branch that does nothing. And the model's
prompt instructions --- which tell it how to interpret \code{episodic}
vs.\ \code{preference} entries --- are stable because the type vocabulary
cannot drift at runtime.

With a well-formed memory block guaranteed by the format layer, the
remaining question is whether to inject it at all. Injecting on every
turn wastes tokens and can introduce stale or irrelevant facts into the
model's reasoning.

%% ------------------------------------------------------------
\section{Four Questions Decide Whether Injection Is Worthwhile}
\label{sec:injection-policy}
%% ------------------------------------------------------------

The problem with unconditional injection is that it treats token cost as
free and model attention as infinite. Neither is true. Every injected
memory block occupies tokens in the \code{HISTORY} slot that could be
used for recent conversational context. If the memories are not relevant
to the current query, the model either ignores them (wasting tokens) or
attends to them incorrectly (producing answers anchored on the wrong
facts). If the memories are already implicit in the recent turns, the
model receives redundant information that increases the risk of
repetitive responses. If the memories are old enough to be misleading,
injection actively harms output quality.

The policy gate encodes four questions as boolean helper functions,
each answerable from data already in the graph state. The four
questions are: Is there at least one memory whose type and content are
relevant to the current user query intent (\code{has\_relevant\_memory})?
Is this memory not already obvious from the recent turns ---
would injection be redundant (\code{is\_redundant\_with\_recent})?
Is the memory recent enough, or explicitly timeless, so that using it
now is safe and not misleading (\code{is\_fresh\_or\_timeless})?
And will injecting this memory improve the model's decision or response
quality more than it increases the token cost
(\code{passes\_cost\_benefit})?

\begin{tipbox}
Configure \code{min\_score} to 0.6--0.7 for most applications. Below
0.5, the relevance filter passes too many marginally related memories
and the cost-benefit question rarely has a ``no'' answer that would have
stopped them. The \code{max\_age\_days} threshold depends on memory
type: set it to 90 for \code{episodic} memories, but leave
\code{preference} and \code{profile} memories as effectively timeless
by passing \code{max\_age\_days=365\,*\,10}.
\end{tipbox}

\begin{lstlisting}[language=Python,
  caption={Four policy helpers and should\_inject\_memory gate},
  label={lst:injection-policy}]
# src/agent/memory/injection_policy.py
from __future__ import annotations
from datetime import datetime, timedelta, timezone
from typing import Callable, Sequence
from langchain_core.messages import BaseMessage

TokenEstimator = Callable[[str], int]

def has_relevant_memory(
    memories: list, query: str, min_score: float = 0.6
) -> bool:
    return any(m["score"] >= min_score for m in memories)

def is_redundant_with_recent(
    messages: list[BaseMessage], memories: list
) -> bool:
    recent_text = " ".join(
        m.content for m in messages[-6:] if m.content
    ).lower()
    for mem in memories:
        if mem["content"][:40].lower() in recent_text:
            return True
    return False

def is_fresh_or_timeless(
    memory: dict, now: datetime, max_age_days: int = 90
) -> bool:
    created = memory["created_at"]
    if created.tzinfo is None:
        created = created.replace(tzinfo=timezone.utc)
    age = (now.replace(tzinfo=timezone.utc) - created).days
    return age <= max_age_days

def passes_cost_benefit(
    memories:         list,
    max_extra_tokens: int,
    estimate_tokens:  TokenEstimator,
) -> bool:
    total = sum(estimate_tokens(m["content"]) for m in memories)
    return total <= max_extra_tokens

def should_inject_memory(
    state:            dict,
    memories:         list,
    now:              datetime,
    max_extra_tokens: int,
    min_score:        float = 0.6,
    max_age_days:     int   = 90,
    estimate_tokens:  TokenEstimator = lambda s: len(s.split()),
) -> bool:
    if not memories:
        return False
    if not has_relevant_memory(memories, _last_query(state), min_score):
        return False
    if is_redundant_with_recent(state.get("messages", []), memories):
        return False
    if not all(is_fresh_or_timeless(m, now, max_age_days) for m in memories):
        memories = [m for m in memories
                    if is_fresh_or_timeless(m, now, max_age_days)]
    if not memories:
        return False
    return passes_cost_benefit(memories, max_extra_tokens, estimate_tokens)

def _last_query(state: dict) -> str:
    from langchain_core.messages import HumanMessage
    for m in reversed(state.get("messages", [])):
        if isinstance(m, HumanMessage):
            return m.content
    return ""
\end{lstlisting}

Three decisions in Listing~\ref{lst:injection-policy} deserve
explanation.

First, \code{is\_redundant\_with\_recent} checks only the last six
messages, not the full history. The full history may contain a distant
mention of the same fact, but a distant mention does not make
re-injection redundant --- the model may no longer be attending to it.
The six-message window reflects the empirical attention boundary for
most instruction-tuned models at production context depths; adjust it if
evaluation data for a specific model suggests a different window.

Second, \code{should\_inject\_memory} does not short-circuit on the
freshness check. Instead, it \emph{filters} stale memories out of the
list and proceeds with the remaining fresh ones. This is deliberate:
a set of memories that is mostly fresh should not be entirely suppressed
because one \code{episodic} entry from eighteen months ago failed the
age check. The freshness filter narrows the set; only an empty set after
filtering produces a \code{False} return on that path.

Third, the default \code{estimate\_tokens} is word-count
(\code{len(s.split())}). This is a deliberate underestimate of actual
token count. The cost-benefit question should err on the side of
injection: if word-count says the block fits, it almost certainly fits
in tokens too. The conservative tokeniser approximation avoids
introducing a \code{tiktoken} dependency in a module that is also used
by lightweight evaluation scripts that do not have the binary available.
Production deployments that need exact token accounting should pass
\code{estimate\_tokens=lambda s: len(get\_encoder().encode(s))} at
construction time.

Chapter~8 introduces mem0 as the persistence and retrieval layer for
long-term memories. The \code{retrieve\_raw\_memories} function built
in \S\ref{sec:injection-pipeline} is the exact integration point: mem0's
similarity-search results arrive as additional \code{RetrievedMemory}
entries appended to the same list, flow through the same six steps, and
are gated by the same four-question policy --- with no modification to
any of the code written in this chapter.
